% Fragmento público generado desde propietarios íntegros.
% Residencia: 52.9.
% Fuente: ../incoming/radion_monografia_20260827/manuscrito/sections/10_minkowski_hodge_lorentz.tex:1-55
\noindent La holonomía orientada admite realizaciones causales y geométricas posteriores en Minkowski, Hodge, Lorentz y Cartan--Holst.\par\medskip

\paragraph{Construcción de la métrica lorentziana desde el bloque APP}

\paragraph{Matriz APP de parámetros \(7\) y \(2\)}

La matriz central
\begin{equation}
B_c=\begin{pmatrix}7&2\\2&7\end{pmatrix}
=9P_++5P_-
\label{eq:Bc}
\end{equation}
se normaliza por su determinante:
\begin{equation}
M_c=\frac{B_c}{\sqrt{45}}
=\exp\left(\frac12\log\frac95\,R\right)
\in SO^+(1,1).
\label{eq:Mc}
\end{equation}
La rapidez interna es
\[
\eta_c=\frac12\log\frac95.
\]
La publicación proyectiva asociada satisface
\begin{equation}
P_B^n=P_++\left(\frac59\right)^nP_-.
\label{eq:PB-contract}
\end{equation}
La misma razón \(5/9\) que aparece como contracción espectral se reconoce aquí
como separación de un modo estable y otro contractivo. La identidad es
operatoria; no convierte todas las apariciones de \(5/9\) en el mismo mapa.

\paragraph{Incidencia \(6\to8\) y cociente de firma causal}

Sea \(M\) la incidencia tipada entre seis caras y ocho vértices. La relación
espectral recuperada es
\begin{equation}
MM^{\mathsf T}=12P_u+4P_-,
\label{eq:MMt}
\end{equation}
donde \(P_u\) proyecta sobre el modo uniforme y \(P_-\) sobre el subespacio
orientado. El cociente superviviente de dimensión cuatro se descompone como
\begin{equation}
Q_4\simeq\R u\oplus E_-.
\label{eq:radion-Q4}
\end{equation}
Declarando negativo el modo uniforme y positivos los tres modos orientados,
se obtiene
\begin{equation}
B_{\mathrm{caus}}=\operatorname{diag}(-1,1,1,1).
\label{eq:minkowski}
\end{equation}
La pantalla de Minkowski no selecciona el estado HMT. Es la realización
causal posterior del cociente que la incidencia ya construyó.

